%! Composition of Functions — Unit 1, Lesson 1.5 — Homework (ANSWER KEY)
\documentclass[10pt]{article}
\usepackage{precalculus-article}
\usepackage{precalculus-key}
\newcommand{\TallMath}[1]{$\displaystyle #1\rule[-1.4em]{0pt}{3.2em}$}

\begin{document}

\pageheader{Unit 1, Lesson 1.5}{Homework --- Answer Key}

\vspace{0.12in}

\begin{enumerate}[itemsep=10pt]

\item The table gives every value of two functions $f$ and $g$. Name the
      handoff before you look anything up.

\smallskip
\renewcommand{\arraystretch}{1.4}
\begin{center}
\begin{tabular}{|c|c|c|c|c|c|}
\hline
$x$ & 1 & 2 & 3 & 4 & 5 \\
\hline
$f(x)$ & 4 & 1 & 5 & 2 & 3 \\
\hline
$g(x)$ & 3 & 5 & 1 & 4 & 9 \\
\hline
\end{tabular}
\end{center}

\begin{enumerate}[label=(\alph*), itemsep=6pt]
  \item $f(g(1))$: \; handoff $= $ \ans{3}, \;
        so $f(g(1)) = $ \ans{5}
  \item $g(f(1))$: \; handoff $= $ \ans{4}, \;
        so $g(f(1)) = $ \ans{4}
  \item $f(g(5))$ has no value. In one sentence, say what goes wrong with the
        handoff.

        \vspace{0.05in}
        \ansline{The handoff is $g(5) = 9$, and 9 is not an input $f$ accepts --- the table has no $f(9)$.}

  \item Complete both rows.

        \smallskip
        \renewcommand{\arraystretch}{1.4}
        \begin{center}
        \begin{tabular}{|c|c|c|c|}
        \hline
        $x$ & 2 & 3 & 4 \\
        \hline
        handoff: $g(x)$ & \ans{5} & \ans{1} & \ans{4} \\
        \hline
        $(f \circ g)(x)$ & \ans{3} & \ans{4} & \ans{2} \\
        \hline
        \end{tabular}
        \end{center}
\end{enumerate}

\item Let $f(x) = x^2 + 2$ and $g(x) = x - 3$. Build both composites and
      simplify.

\begin{work}
  f(g(x)) &= f(x - 3) \\
          &= (x - 3)^2 + 2 \\
          &= x^2 - 6x + 9 + 2 \\
          &= x^2 - 6x + 11 \\
  g(f(x)) &= g(x^2 + 2) \\
          &= (x^2 + 2) - 3 \\
          &= x^2 - 1
\end{work}

\begin{enumerate}[label=(\alph*), itemsep=5pt]
  \item $(f \circ g)(x) = $ \ans{$x^2 - 6x + 11$}
  \item $(g \circ f)(x) = $ \ans{$x^2 - 1$}
  \item One input settles whether they are the same function. At $x = 0$:
        $f(g(0)) = $ \ans{11} and $g(f(0)) = $ \ans{$-1$}.
\end{enumerate}

\item A \$45 backpack, the same \$10 coupon, and the same 6\% tax:
      $c(p) = p - 10$ and $t(x) = 1.06x$.

\begin{work}
  t(c(45)) &= t(35) \\
           &= 1.06(35) \\
           &= 37.10 \\
  c(t(45)) &= c(47.70) \\
           &= 47.70 - 10 \\
           &= 37.70
\end{work}

\begin{enumerate}[label=(\alph*), itemsep=5pt]
  \item Coupon first: handoff \$\ans{35}, total \$\ans{37.10}.
  \item Tax first: handoff \$\ans{47.70}, total \$\ans{37.70}.
  \item The gap is \$\ans{0.60} --- the same gap as the \$60 jacket. Say
        in one sentence why the price does not affect it.

        \vspace{0.05in}
        \ansline{Taxing first charges 6\% of the \$10 coupon, and $0.06 \times 10 = 0.60$ no matter the price.}
\end{enumerate}

\item Take each function apart, then put one back together.

\begin{enumerate}[label=(\alph*), itemsep=6pt]
  \item $h(x) = \sqrt{4x + 1}$: \quad inner $g(x) = $ \ans{$4x + 1$},
        \quad outer $f(u) = $ \ans{$\sqrt{u}$}
  \item $k(x) = 3(x - 2)^5$: \quad inner $g(x) = $ \ans{$x - 2$},
        \quad outer $f(u) = $ \ans{$3u^5$}
  \item In Lesson 1.4 you wrote the horizontal shift $g(x) = f(x + 6)$.
        Write it as a composition: $g(x) = f(h(x))$ where
        $h(x) = $ \ans{$x + 6$}, and the handoff is \ans{$x+6$}.

        Because the handoff is what $f$ gets \ans{asked}, the graph moves
        6 units to the \ans{left}.
\end{enumerate}

\item Let $f(x) = x^2$ and $g(x) = 2x - 1$. Which of the following is
      $(f \circ g)(x)$?

\begin{enumerate}[label=\Alph*., itemsep=3pt]
  \item $2x^2 - 1$
  \item \ans{$4x^2 - 4x + 1$}
  \item $4x^2 + 1$
  \item $2x^3 - x^2$
\end{enumerate}

\end{enumerate}

\vspace{0.1in}

\boxguard
\begin{extensionbox}
Let $f(x) = 5x - 2$ and let $I(x) = x$ be the identity function.

\begin{enumerate}[label=(\alph*), itemsep=4pt]
  \item Show that $f(I(x))$ and $I(f(x))$ both come out to $f(x)$ --- with
        $I$ in the chain, the handoff is the number itself.
  \item Now hunt for a different function $g$ so that $f(g(x)) = x$: a partner
        whose handoff comes back out as the number you started with.
\end{enumerate}

\vspace{0.05in}
\ansline{(a) $f(I(x)) = f(x) = 5x-2$ and $I(f(x)) = I(5x-2) = 5x-2$ --- $I$ changes nothing either way.}
\ansline{(b) $g(x) = (x+2)/5$, since $f\big((x+2)/5\big) = 5 \cdot (x+2)/5 - 2 = x$. That is $f^{-1}$.}
\end{extensionbox}

\vspace{0.1in}

\boxguard[9]
\begin{notesbox}{Preview of Lesson 1.6}
The extension asks for a function that \textit{undoes} $f$. Lesson 1.6,
\textit{Inverse Functions}, gives that partner a name and a notation ---
and the test for it is a composition: $f\big(f^{-1}(x)\big) = x$, the identity
function from Part 5.
\end{notesbox}

\end{document}
